\documentclass[11pt]{amsart}
\usepackage[T1]{fontenc}
\usepackage{amsmath,amssymb,amsthm}
\usepackage{graphicx}
\usepackage[colorlinks=true,linkcolor=blue,citecolor=blue,urlcolor=blue]{hyperref}

\newtheorem{theorem}{Theorem}
\newtheorem{remark}{Remark}

\begin{document}

\title{An \(\ell_1\)-Sum Obstruction to a Directional Smoothness Property}
\author{Run fa\_banach\_001}
\date{2026-06-16}
\maketitle

\section*{Status}

This is a candidate partial result for the question in Remark 2.11 of
Dilworth--Garrig\'os--Hern\'andez--Kutzarova--Temlyakov
\cite{DGHKT}. It does not solve the full question. It rules out a natural
class of non-uniformly smooth spaces as possible examples.

\section*{Source Question}

In Definition 1.12 of \cite{DGHKT}, a pair \((X,\mathcal D)\) has property
\(\mathcal D(s,c_1)\) if, for some \(s>1\) and \(c_1>0\),
\[
  \operatorname{dist}(f,[\varphi])
  \leq \|f\|\bigl(1-c_1 |F_f(\varphi)|^s\bigr)
\]
for every \(f\neq 0\) and every atom \(\varphi\in\mathcal D\), where \(F_f\)
is an associated norming functional. Remark 2.11 asks whether this property
could hold for some non-uniformly smooth Banach space and some dictionary.

\begin{center}
\includegraphics[width=\textwidth]{figures/open_problem_crop.png}
\end{center}

\section*{Proof Intuition}

The obstruction is that an \(\ell_1\) direct sum has two independent norm
components. Given any proposed unit atom \(\varphi=(u,v)\), one can build
vectors \(x_t\) whose two components point in exactly the same directions as
\(u\) and \(v\), but with very different relative sizes. Every norming
functional for \(x_t\) still evaluates \(\varphi\) with modulus one. However,
because one component of \(x_t\) is much larger than the matching component of
the atom, subtracting a scalar multiple of \(\varphi\) cannot reduce the norm
by a fixed proportion. The best distance ratio tends to one, contradicting
any inequality with a fixed positive \(c_1\).

\section*{Theorem}

Let \(E\) and \(F\) be nonzero Banach spaces over \(\mathbb K=\mathbb R\) or
\(\mathbb C\), and let
\[
  X = E\oplus_1 F,\qquad \|(e,f)\|=\|e\|+\|f\|.
\]
Then for every unit vector \(\varphi\in X\), every \(s>1\), and every \(c>0\),
the inequality
\[
  \operatorname{dist}(x,[\varphi])
  \leq \|x\|\bigl(1-c |G_x(\varphi)|^s\bigr)
\]
fails for some \(x\neq 0\), no matter how one chooses associated norming
functionals \(G_x\). Consequently no nonempty dictionary in a nontrivial
\(\ell_1\) sum can satisfy property \(\mathcal D(s,c)\).

\section*{Proof}

Let \(\varphi=(u,v)\in X\) with \(\|\varphi\|=1\). Put
\[
  A=\|u\|,\qquad B=\|v\|,
\]
so \(A+B=1\).

First suppose \(A,B>0\). Set \(e=u/A\) and \(y=v/B\), so
\(\|e\|=\|y\|=1\). If \(A\leq B\), take
\[
  x_t=(e,t y),\qquad t>0.
\]
Then \(\|x_t\|=1+t\). Any associated norming functional for \(x_t\) has the
form \(G_t=(e_t^*,y_t^*)\), with \(\|e_t^*\|=\|y_t^*\|=1\),
\(e_t^*(e)=1\), and \(y_t^*(y)=1\). Hence
\[
  G_t(\varphi)=e_t^*(u)+y_t^*(v)=A+B=1.
\]
For any scalar \(\lambda\),
\[
  \|x_t-\lambda\varphi\|
  = |1-\lambda A| + |t-\lambda B|.
\]
Since \(A/B\leq 1\),
\[
  |1-\lambda A|+|t-\lambda B|
  \geq |1-\lambda A|+\frac AB |t-\lambda B|
  \geq 1-\frac AB t.
\]
Equality is attained by taking \(\lambda=t/B\). Therefore
\[
  \operatorname{dist}(x_t,[\varphi])=1-\frac AB t.
\]
Dividing by \(\|x_t\|=1+t\) and letting \(t\downarrow 0\) gives
\[
  \frac{\operatorname{dist}(x_t,[\varphi])}{\|x_t\|}\longrightarrow 1,
  \qquad |G_t(\varphi)|=1.
\]
Thus no fixed \(c>0\) can make the desired inequality true for all \(t\).

If \(A>B\), the same argument with \(x_t=(t e,y)\) gives
\[
  \operatorname{dist}(x_t,[\varphi])=1-\frac BA t,\qquad
  \|x_t\|=1+t,\qquad |G_t(\varphi)|=1,
\]
again contradicting any fixed positive \(c\).

It remains to consider the pure cases. If \(A=0\), then \(B=1\). Choose any
unit vector \(e\in E\), put \(y=v\), and set \(x_t=(e,t y)\). Every norming
functional for \(x_t\) evaluates \(\varphi\) as \(1\), while
\[
  \operatorname{dist}(x_t,[\varphi])=1,\qquad \|x_t\|=1+t.
\]
Letting \(t\downarrow0\) again contradicts any positive \(c\). The case
\(B=0\) is symmetric. This proves the atom-level obstruction for every unit
\(\varphi\), and therefore for every nonempty dictionary.
\qed

\section*{Limitations And Novelty Check}

The theorem is only a partial negative result. It does not prove that property
\(\mathcal D(s,c_1)\) forces uniform smoothness, and it does not rule out all
non-uniformly smooth norms or all possible thin dictionaries. It shows that
non-uniform smoothness coming from a nontrivial \(\ell_1\) direct-sum corner is
incompatible with the property at the level of every single atom.

The run indexes were searched for arXiv:1909.13536, the paper title,
\(\mathcal D(s,c_1)\), non-uniformly smooth Banach spaces, and WCGA
Lebesgue-type terms. A web search found the source arXiv page but no later
exact answer to Remark 2.11; a Semantic Scholar API citation query was
rate-limited. Novelty confidence is therefore modest and should be checked by
a human reviewer.

\begin{thebibliography}{9}

\bibitem{DGHKT}
S. Dilworth, G. Garrig\'os, E. Hern\'andez, D. Kutzarova, and V. Temlyakov,
\emph{Lebesgue-type inequalities in greedy approximation},
arXiv:1909.13536.

\end{thebibliography}

\end{document}
